% ------------------------------------------------------------------
%  Methods.  Journal prose (Nature Portfolio order).  Provenance of every
%  statement -- code paths, configuration files, slide sources -- is kept
%  in notes/sources.md, not in the manuscript.
% ------------------------------------------------------------------

\subsection*{Ethics and study oversight}
This retrospective study analysed existing clinical records and \mngs reports. Patients were not contacted, and pipeline outputs were not used in clinical care. The study was approved by the institutional review boards of Kaohsiung Medical University Hospital (\kmuh; \todo{protocol number}), Taipei Veterans General Hospital (\tvgh; \todo{protocol number}) and Tri-Service General Hospital (\tsgh; \todo{protocol number}); each board waived the requirement for informed consent because the study used retrospectively collected, de-identified data.

Direct identifiers were removed before records were submitted to a commercial language-model service. Submissions were made through the provider's application programming interface under terms that exclude the use of submitted data for model training, and reference labels were withheld from all model inputs. \toolname was evaluated as research software and returned candidate pathogens for physician review rather than a clinical diagnosis.

\subsection*{Study design and participants}
We conducted a retrospective, multicentre diagnostic-accuracy study at three tertiary referral centres in Taiwan. Participants underwent \mngs for suspected lower respiratory tract infection at \kmuh, \tvgh or \tsgh between August 2023 and September 2025. The index test was \toolname, and the reference standard was the causative-pathogen list adjudicated by the treating physicians. Comparators were conventional microbiological read-outs interpreted at face value and four general-purpose language models prompted directly with the same patient data.

The three hospitals maintained a shared registry of 105 \mngs specimens (\kmuh, 51; \tvgh, 35; \tsgh, 19), together with demographics, comorbidities, conventional microbiology, \mngs reports, treatment changes and outcomes (Table~\ref{tab:cohort}). Two evaluation cohorts were defined before model inference (Extended Data Fig.~\ref{fig:ed1}). The primary cohort included 55 patients (\kmuh, 30; \tvgh and \tsgh, 25) with a complete standardized record, a per-specimen \mngs report and at least one adjudicated causative pathogen at a pulmonary site. Patients whose adjudicated final diagnosis was infection at another site, such as urinary tract infection, or no infection, were therefore not evaluated, even where \mngs had been sent for suspected pneumonia. This restricts the estimates reported here to patients in whom lower respiratory tract infection was ultimately confirmed, a population that cannot be identified at the time the test is ordered; the consequences for the interpretation of precision are considered in the Discussion. \todo{Give the number excluded on this ground and the site of infection in each, and state whether the exclusion was applied to the registry or to the delivered set.} The direct-prompting cohort included the 30 \kmuh patients and the 11 \tvgh and \tsgh patients whose records were complete in every section ($n=41$). \todo{Reasons for the two-hospital exclusions (54 registry candidates, 48 with complete records, 25 in the primary cohort).} An earlier series of 35 \kmuh intensive-care patients, analysed with the first pipeline version, was used only for the organism-burden comparison (Fig.~\ref{fig:benchmark}c,d) \todo{state its overlap with the primary cohort}.

\subsection*{Metagenomic sequencing and conventional microbiology}
\mngs was ordered by the treating physicians as part of clinical care and performed by one clinical laboratory (Asia Pathogenomics, New Taipei City, Taiwan) using a previously described workflow\cite{chou2025investigating,takahashi2026diagnostic}. Specimens were predominantly bronchoalveolar lavage fluid (60\%) and blood (23\%; Table~\ref{tab:cohort}). Briefly, specimens were heat-inactivated and mechanically lysed. Nucleic acids were extracted without host depletion using the TIANMicrobe Pathogen DNA Kit (Tiangen) and, for RNA, the QIAamp Viral RNA Mini Kit (Qiagen)\cite{chou2025investigating}. Libraries were prepared with unique dual-index adapters (TIANSeq Direct Fast DNA Library Prep Kit, Tiangen) and sequenced as 75-bp single-end reads on an Element Biosciences AVITI system, yielding approximately 40 million reads per sample\cite{takahashi2026diagnostic}.

Reads were quality-filtered with Trimmomatic\cite{bolger2014trimmomatic} and PRINSEQ\cite{schmieder2011quality}. Human reads were removed by alignment to T2T-CHM13\cite{nurk2022complete} with BWA-MEM\cite{li2013aligning}, and the remaining reads were aligned to a curated database containing 5,884 viral, 18,432 bacterial, 3,146 fungal and 370 parasitic genomes\cite{takahashi2026diagnostic}. Extraction, library and no-template negative controls were included in each run. Organisms were reported when normalized abundance exceeded ten times the highest value in a concurrent negative control. For fastidious, low-biomass pathogens such as \textit{Mycobacterium tuberculosis}, uniquely mapping reads were required in the specimen and absent from all controls\cite{takahashi2026diagnostic}. \toolname used only the resulting clinical report---the reported organisms, read counts and ranks in the negative-control comparison---and did not process raw sequencing reads.

Conventional microbiology was performed in each hospital's accredited clinical laboratory as part of routine care: culture with identification by mass spectrometry or automated biochemical systems; the BioFire \filmarray Pneumonia Panel (bioM\'erieux) on lower respiratory specimens, which reports typical bacteria semi-quantitatively and atypical bacteria, viruses and resistance genes qualitatively\cite{buchan2020practical,chou2025investigating}; the Platelia \textit{Aspergillus} galactomannan enzyme immunoassay (Bio-Rad) in serum and lavage fluid\cite{hung2021investigating}; and targeted molecular tests. Results were taken as reported. Within the pipeline, a \filmarray bacterial target contributes evidence in proportion to its semi-quantitative bin, and galactomannan is considered positive at an index of $\ge$1.0 in lavage fluid and $\ge$0.5 in serum, in line with consensus definitions\cite{donnelly2020revision}.

\subsection*{Reference standard}
The reference standard was established by a clinical review committee, following the approach used in clinical mNGS studies in which physicians adjudicate the clinical significance of sequencing findings\cite{wilson2019clinical,miller2019laboratory}. The committee comprised three infectious-diseases physicians from the participating medical centres. For each patient, the committee reviewed the presentation, imaging, blood counts, inflammatory markers, conventional microbiology, \mngs report, clinical course and response to treatment. It classified each organism detected by any method as a causative pathogen of the current episode or as a colonizer, contaminant, reactivated virus or bystander. Clinical context was used to adjudicate organisms that culture cannot confirm, including viruses, fungi and fastidious bacteria. The resulting causative pathogen list (one to eight organisms per patient) was recorded in the registry. Two label revisions affecting four patients were finalized before evaluation, and the final labels were used for all reported analyses. Labels were available only to the evaluation scripts and were excluded from pipeline inputs, prompts and configuration. Because the committee reviewed the \mngs report, the reference standard was not independent of the index test. The absence of independent adjudication and blinding during rule development is addressed in the Discussion.

\subsection*{Data extraction and standardization}
Source records comprised per-patient workbooks at \kmuh and a shared master workbook with case-presentation slides at \tvgh and \tsgh. Rule-based parsers divided workbook contents into clinical sections and extracted slide text; information contained only in images was transcribed from rendered slides. Slides were linked to registry cases using patient, date and file metadata.

A language model (GPT-5.6 Luna, OpenAI) then converted each raw section to a fixed schema for nine clinical domains: admission diagnosis, underlying conditions, complete blood count, other laboratory values, cultures, \filmarray, galactomannan, molecular microbiology and imaging. The \mngs report was retained as a tenth input. The same prompt and output schema were used for all records. Clinical observations were restricted to a window anchored on collection of the \mngs specimen, and records without a parsable time were excluded. Hospital-specific patient identifiers and specimen codes were used to link records while preventing collisions between institutions.

\subsection*{Overview of \toolname}
\toolname analyses one patient at a time in five stages (Fig.~\ref{fig:overview}). First, modality-specific language models convert the standardized record into structured clinical evidence. Second, prespecified rules combine the clinical evidence with the \mngs signal to make an initial pathogen selection. Third, a safety review identifies unselected organisms that warrant additional evaluation but cannot change the initial selection. Fourth, the system evaluates the retained candidates against published cases, sequencing strength, direct microbiological evidence and anatomical plausibility, after which a rule-based integration step fixes the final list. Finally, a language model explains the fixed decision.

Language models were therefore used to extract or assess evidence, whereas version-controlled rules made every selection. Each model task used a fixed prompt and structured output schema (Supplementary Table~5). The pipeline recorded the rule applied at each decision and cryptographic hashes of the corresponding inputs. Its output comprised the selected pathogen list and, for every organism considered, the evidence for and against selection and the rule that determined its disposition. Explanation generation occurred only after selection and could not alter the list.

\subsection*{Modality-specific evidence extraction}
Five language-model agents (GPT-5, OpenAI) processed distinct parts of the standardized record: host factors and laboratory measurements; culture; \filmarray and galactomannan; molecular microbiology; and imaging. Each agent received only its assigned clinical domain, used a fixed prompt and output schema and had no access to the other agents or the reference labels. The host agent summarized immunosuppression, malignancy, corticosteroid exposure, renal failure and related factors as prespecified vulnerability and opportunistic-infection categories. The imaging agent characterized the likelihood and anatomical pattern of infection and extracted organism-specific radiological clues. Imaging clues were retained as clinical context but were not treated as direct organism confirmation, which required culture, \filmarray, galactomannan or molecular evidence.

Each organism in the \mngs report was assigned to its source specimen and anatomical site. A rule-based summary then represented, for each organism, the evidence from each modality (support, no support or unknown), whether the specimen was anatomically aligned with the suspected infection, and the relevant host context. Severe host vulnerability broadened consideration of otherwise borderline organisms but could not by itself cause selection (Supplementary Tables~1 and 2).

\subsection*{Characterization of \mngs signals}
Each organism in the \mngs report was characterized by its read count, its relative abundance among organisms detected in the same run, its rank against negative controls, its dominance within the specimen and the anatomical relevance of that specimen (Supplementary Table~1). Repeated detections in one patient were resolved by a prespecified hierarchy based on negative-control rank and read burden. Organism names were normalized with a fixed dictionary that merged synonyms (for example, CMV and human cytomegalovirus) but did not collapse distinct species.

\subsection*{Initial pathogen selection}
The initial scorer combined the \mngs features with specimen alignment, direct hospital-test support and host context. Each organism was assigned to one of five ordered causative levels, from strongly supported to background (Supplementary Table~2). Strong support required a favourable negative-control rank together with high abundance or within-specimen dominance. Intermediate signals could be strengthened by exact-organism evidence from another microbiological test or by prespecified rules for clinically important pathogens; weak or anatomically discordant signals were downgraded.

Before selection, prespecified guardrails reduced common sources of overinterpretation, including respiratory \textit{Candida} without invasive evidence, oral and aspiration flora, skin and environmental organisms, reactivated or low-specificity viruses and weak non-dominant detections (Supplementary Table~3). Strongly supported organisms were selected directly. Borderline organisms required compatible clinical context, direct microbiological support or a pathogen-specific rule; host vulnerability alone was insufficient. Respiratory \textit{Candida} and generic yeasts always required evidence of invasive infection.

When no organism met the selection criteria, the strongest candidate that passed the exclusion rules was reported as the best available. The system could also consider an organism absent from the \mngs report when an exact-species culture, \filmarray or galactomannan result, or targeted molecular assay provided direct evidence from a clinically relevant specimen. Selection in this route depended on the hospital test rather than sequencing features. IgM serology alone was not treated as direct evidence; \textit{Chlamydophila} or \textit{Mycoplasma pneumoniae} supported only by IgM was retained for contextual review.

\subsection*{Safety review of unselected organisms}
A language-model safety review (GPT-5.6 Luna) examined organisms not selected by the initial rules. It received the deterministic result, the structured hospital evidence and specimen assignments, but could neither remove nor reorder an initial selection. Its prompt emphasized that immunocompromise or membership in a clinically important pathogen group was insufficient for selection on its own and restated the prespecified rules for colonization, reactivation, oral flora, environmental organisms and \textit{Candida} (Supplementary Note~1).

The review separated organisms warranting further evidence assessment from those with low specificity or insufficient support. A rule-based check then standardized organism names and removed candidates supported only by anatomically irrelevant specimens or weak clinical context. Thus, the review could expand the set undergoing evidence assessment but could not promote an organism directly to the final pathogen list.

\subsection*{Literature evidence and patient-level case matching}
For every initial selection and additional candidate retained by the safety review, PubMed was searched through NCBI E-utilities using the organism name and lower-respiratory-tract infection terms. Searches included publications available by 30 June 2026 and returned up to ten articles per organism. A language model (GPT-5.6 Luna) assessed each title and abstract for human clinical evidence that the exact species causes infection at the relevant anatomical site. Supporting evidence required an exact-species match, human infection at the target site and a traceable excerpt from the source. When too few articles could be assessed, the literature analysis abstained rather than treating limited evidence as evidence against infection (module A; Supplementary Table~4).

The same articles were then compared with a label-blind summary of the individual patient (module A1). This summary included specimen type and site, clinical syndrome, host factors, imaging and laboratory phenotype and timing, but excluded the reference label, initial selections, read counts, culture and PCR results and all previous review rationales. The model extracted the reported disease context and graded agreement in anatomical site, syndrome, specimen, host, phenotype and timing. Deterministic rules converted these structured observations to strong match, partial match, mismatch or insufficient evidence. Exact-species causal evidence with compatible site and syndrome was required for a positive match; explicit colonization, a conflicting site or syndrome, non-human evidence or broader taxonomic evidence could not support rescue. Only articles with traceable supporting or opposing excerpts entered the aggregate case-match score.

\subsection*{Independent evidence dimensions}
Three additional dimensions summarized evidence independent of article case matching (Supplementary Table~4). Sequencing strength captured both absolute burden and relative prominence within the specimen (module B). Direct microbiological evidence was graded by assay specificity and anatomical relevance, with exact-species detection in a sterile site or tissue receiving the strongest grade (module C). Results without a usable specimen site or with contamination concerns were treated as context rather than confirmation, and respiratory \textit{Candida} culture alone was not considered evidence of invasive disease. A final check recorded anatomical coherence and colonization concerns for audit (module D). It did not determine selection, but an explicit anatomical mismatch could not be overridden.

\subsection*{Final evidence integration and pathogen selection}
The final rule (\rfive) retained all organisms selected initially and added a previously unselected organism through either of two rescue routes. The first required exact-species evidence from a sterile site or tissue. The second required a prespecified level of concordance with published cases, anatomical alignment with the suspected infection and either high absolute sequencing burden or strong relative prominence within the specimen. The complete rule and thresholds are reported in Supplementary Table~4.

Initial selections retained their original order, and rescued organisms were appended without an inferred clinical ranking. For each patient, the pipeline recorded the selection route for every organism and the reason each remaining candidate was not selected. Failure to satisfy a rescue route was reported as insufficient evidence for selection, not as evidence that infection was absent.

\subsection*{Explanation generation, audit and delivery}
After final selection, candidate-specific microbiology within two days before or after collection of the \mngs specimen was assembled for explanation; the complete time course was retained separately. A language model (GPT-5.6 Luna) generated one explanation for each selected organism from these evidence views and the fixed final decision. The generator was instructed that the selected organisms and their order were immutable. Automated checks required exactly one entry per selected organism, prohibited any organism outside the selected list and verified the presence of all required fields; a failed check blocked delivery.

The delivered record stated the disposition of every organism detected by any method. For selected organisms, it reported the selection route, supporting and opposing findings, relevant literature and sequencing evidence, and specimen site and timing. For non-selected organisms, it reported the applicable exclusion rule or unmet evidence criterion and clarified that non-selection did not exclude infection. Each statement was linked to the corresponding structured evidence. Outputs were delivered as a summary table and an auditable record (Fig.~\ref{fig:explain}) and were labelled as research results pending clinical review.

\subsection*{Conventional read-out comparators}
Each conventional read-out was scored at face value by treating its reported organisms as predictions: every organism on the \mngs report; every \filmarray or galactomannan target reported positive; and every organism isolated in culture within the evidence window. \todo{Confirm these definitions and the window against the scoring script, and state the handling of normal flora, no growth and pending results.} No read-out was reinterpreted or filtered. These comparators represent the information a physician receives before any integration.

\subsection*{Direct-prompting comparison}
To determine whether performance arose from the structured pipeline rather than from access to a language model alone, four general-purpose models were evaluated by direct prompting in 41 patients. Each model received a single request containing the same nine standardized clinical sections and the laboratory \mngs report, including organism names, read counts, negative-control ranks and specimen types. Direct identifiers and all fields derived from the reference standard or from prior pipeline decisions were removed.

The Traditional Chinese prompt asked the model to act as an infectious-diseases reviewer and emphasized that organism detection does not establish causation. It specified common concerns involving colonization, anatomical site, \textit{Candida}, environmental organisms and viral reactivation; allowed zero or more pathogens; required species-level predictions; and requested supporting evidence, opposing evidence, excluded organisms and limitations in a structured response (Supplementary Note~2). \gptsol, \gptluna, \gptterra and \gptastra (OpenAI) were queried through the Responses API with medium reasoning effort, a JSON output schema and a maximum of 4,000 output tokens. Tools, external retrieval and conversation history were disabled. Each model was run once per patient, and the response, prompt hash and payload hash were stored before scoring.

\subsection*{Outcome measures and scoring}
The unit of evaluation was the (patient, organism) pair. Predicted and reference names were normalized using a fixed synonym table (for example, CMV to human cytomegalovirus and PJP to \textit{P.~jirovecii}) and matched at exact-species level. Genus-level matches were not accepted; thus, a prediction of \textit{Candida} against a reference of \textit{Candida tropicalis} contributed one false positive and one false negative.

The primary outcomes were precision, recall and F1, micro-averaged across all patient--organism pairs. Secondary outcomes were the corresponding patient-level macro-averages; the proportions of patients in whom at least one, all or exactly the reference set of pathogens was identified; and the number of patients for whom a method returned no organism (Supplementary Note~5). Both evaluation cohorts required at least one adjudicated pathogen, so precision does not capture false-positive predictions in pathogen-negative patients. Explanations were tested for structural completeness and consistency with the fixed decision, but were not independently rated for clinical usefulness or clause-level faithfulness to the source record; \todo{a reader study rating the delivered explanations, and the effect of the explanation on antimicrobial decisions, remain to be done}.

\subsection*{Secondary registry analyses}
The following analyses were exploratory and were not used to develop or evaluate the pathogen-selection rules.

\subsubsection*{Concordance of sequencing and conventional microbiology}
To characterize the information provided by the \mngs report before multimodal integration, we analysed registry cases with an adjudicated pathogen list in the shared workbook ($n=33$; Supplementary Note~6). \kmuh cases were not included in this analysis because their adjudications were stored separately in per-patient workbooks. For each case, we compared the organisms and read counts in the \mngs report with organisms detected by culture and non-culture tests and with the adjudicated pathogen list. Names were normalized and matched at exact-species level using the same rules as the primary evaluation. Entries explicitly indicating that no pathogen was present were excluded.

An \mngs detection was considered corroborated when the same species was identified by any conventional test in the same patient. For adjudicated pathogens absent from the \mngs report, we recorded the detecting assay, specimen type and collection time relative to receipt of the \mngs specimen. Conventional assays included culture, PCR and multiplex panels, viral-load testing, antigen testing and IgM or antibody serology. Bronchoalveolar lavage fluid and sputum were classified as non-sterile specimens; blood, cerebrospinal fluid, tissue, ascites, pus and abscess material were classified as sterile. The ability of read count alone to distinguish causative from non-causative detections was summarized by the area under the receiver operating characteristic curve, computed for the 80 of 104 detections with a recorded read count.

\subsubsection*{Composition and timing of the clinical record}
To quantify the evidence available for physician interpretation (Fig.~\ref{fig:explain}a--d), we analysed the complete per-admission source workbooks from \kmuh. These records contained underlying conditions; admission, transfer and discharge diagnoses; culture reports with collection and reporting times; full-text radiology reports; all syndromic-panel targets; galactomannan measurements; and blood count and chemistry results. Unlike the registry summary, they retained negative results, timestamps and the original report wording.

Of 36 workbooks, 32 distinct admissions were analysed. One lacked microbiology, imaging and panel data; one was labelled as a duplicate; and two were byte-identical copies of another admission that had undergone repeat \mngs. Repeat sequencing of the same admission was counted once. We defined a result as one culture report, radiology study, panel target, galactomannan measurement or recorded analyte value. Organisms were counted once per admission and read-out after synonym normalization; galactomannan was considered to identify \textit{Aspergillus} at an index of $\geq0.5$ in serum or $\geq1.0$ in lavage, and resistance markers were not counted as organisms. Sterile-site cultures comprised blood, ascites, pleural fluid, tissue, cerebrospinal fluid, abscess, deep-wound and intravascular-catheter specimens. Culture turnaround was the interval from specimen collection to the final report.

Radiology wording was classified as qualified when a report used terms such as ``probably'', ``suspect'', ``favour'', ``cannot be excluded'', ``questionable'', ``equivocal'', ``may represent'' or ``suggestive of''. Formulaic recommendations for clinical correlation or differential diagnosis were not considered qualifications. Only aggregate results from these protected source records were exported for analysis.

For each read-out, we summarized the interval between specimen collection and receipt of the \mngs specimen, the interval between admission and \mngs collection and the proportion of patients in whom that read-out identified an organism. Collection intervals within 90 days were summarized by the median, interquartile range and 5th--95th centiles; larger intervals were treated as date-entry errors. Because the registry records conventional tests almost exclusively when an organism was detected (225 of 226 dated culture entries were positive), these distributions describe the timing of organism-identifying evidence, not the complete testing burden. They therefore cannot estimate culture yield or the effect of previous antimicrobial treatment.

The registry recorded arrival and reporting dates for \mngs specimens, but these dates were identical for more than half of the records and were not used to estimate sequencing turnaround. These analyses were descriptive, and no inferential test was applied. Because adjudicators had access to the \mngs report, comparisons with the reference standard share the same non-independence as the primary evaluation. Only de-identified per-detection data and aggregate counts were exported from the protected registry.

\subsubsection*{Physiological trajectory around the sequencing report}
% tables/make_trajectory.py; reads the registry workbook and the KMUH whole-record extracts (both PHI, outside the repository) and writes only aggregates
We descriptively examined physiological trajectories around \mngs specimen collection using machine-readable extracts of the complete \kmuh discharge record. Each dated laboratory value and vital sign retained a link to its source page. Extracts were available for 49 of 51 \kmuh registry admissions; 48 also had a dated \mngs specimen and recorded discharge status and were included (31 deaths and 17 survivors).

The collection date was defined as day~0. Measurements were grouped into five windows: days $-14$ to $-1$, 0--2, 3--7, 8--14 and 15--30. For each patient and window, we first calculated the median value of each marker so that patients with frequent measurements contributed once per interval; group summaries were then calculated from these patient-level medians. Pre-collection and day 3--7 values were compared with a two-sided Wilcoxon signed-rank test among patients with measurements in both windows, without adjustment for multiple comparisons. Fever was summarized as the maximum temperature within each window and defined as a value above 38\,\textdegree C. Improvement was defined as movement in the clinically favourable direction between the pre-collection and day 3--7 windows; improvement followed by worsening required a subsequent reversal during days 8--30.

Values outside prespecified physiological ranges were excluded. Temperature values were reconstructed from the source table layout and underwent a documented correction for a systematic transcription error, which was checked against adjacent readings from the same admission.

This analysis was observational and descriptive. All patients underwent \mngs, and there was no unexposed comparator. Antipyretics, corticosteroids, sedation and ventilator settings could directly affect temperature and gas exchange but were not available in a form suitable for adjustment. Antimicrobial start and stop dates were also unavailable; therefore, neither time from reporting to treatment modification nor coverage of the adjudicated pathogen could be determined. The analysis was not used to make causal claims or to propose a monitoring test.

\subsubsection*{Registry treatment and outcome analysis}
% tables/make_impact.py (main:154 impact block, main:236 outcome block); tests implemented in the same
% file (fisher_or, logistic_or) because the analysis environment has no scientific Python stack.
Treatment modification, its direction, its attribution to \mngs, mechanical ventilation and discharge status were extracted for all 105 registry specimens and summarized as counts and percentages. Denominators were reported when attribution or ventilation status was missing. The number of organisms per \mngs report was compared across specimen types and patient strata. Because specimen type was associated with disease severity, stratum comparisons were repeated within bronchoalveolar-lavage specimens, and only these restricted comparisons were interpreted. Organism-level agreement with culture was calculated after the same name normalization used in the primary evaluation and cross-checked against the registry's existing comparison.

Because every registry patient underwent \mngs, the data could not estimate the effect of sequencing relative to no sequencing. We therefore compared patients whose treatment modification was retrospectively attributed by the treating team to \mngs with those whose modification was not. This attribution was neither assigned nor assessed under blinding. In-hospital mortality was compared with Fisher's exact test. Odds ratios and 95\% confidence intervals used the Haldane--Anscombe correction and were additionally estimated with logistic regression adjusted for intensive care at sampling, acute respiratory distress syndrome, active malignancy, corticosteroid treatment and age; Wald intervals are reported.

Comparisons between patients whose treatment changed for any reason and those whose treatment did not were reported descriptively because of evident confounding by indication. The registry did not contain antimicrobial names or administration dates, precluding assessment of treatment appropriateness or time to appropriate therapy. Time-to-event outcomes were also not analysed because patients had to survive long enough to receive a change attributed to a specimen collected after admission. All tests were two-sided, without adjustment for multiple comparisons. These secondary analyses were considered hypothesis-generating, and only aggregate results were exported from the protected registry.

\subsection*{Computational reproducibility}
A single orchestrator validated the input to each stage, executed the stages in a fixed order and recorded all commands and output hashes. It contained no pathogen-selection logic and did not overwrite previous outputs. All reported final decisions were regenerated from archived intermediate files, and direct-prompting metrics were recomputed from stored model responses. The current code release reproduced the selected organisms for every patient in the primary cohort. Two implementation differences from the version used for the reported runs did not alter any selection and are described in Supplementary Note~4. The test suite runs offline without network access or credentials.

\subsection*{Statistics and reproducibility}
No statistical method was used to predetermine sample size; the cohorts included every registry patient who met the eligibility criteria when the datasets were finalized. Patients were not randomized because all methods were evaluated on the same records. Reference labels were withheld from the language-model stages, and each method's predictions were recorded before labels were accessed by the evaluation scripts. However, the investigators who developed the rule-based scoring system had access to the labels, and no patients were held out during rule development. The reported diagnostic performance is therefore a development-set result (Discussion).

Patients were excluded only for missing inputs and before model inference (Extended Data Fig.~\ref{fig:ed1}). The clinical records were not publicly available and were not used to train the evaluated language models. Precision, recall and F1 are reported as point estimates. Rule-based stages were exactly reproducible, whereas each language-model stage was run once per patient and its output was retained; the reported results therefore do not capture run-to-run model variability. \todo{Add patient-level bootstrap 95\% confidence intervals and paired differences in F1 versus each comparator.}

\subsection*{Reporting summary}
Further information on research design is available in the Nature Portfolio Reporting Summary linked to this article.

\subsection*{Data availability}
The patient-level clinical records, \mngs reports, reference labels, model inputs and model outputs contain protected health information and cannot be shared publicly. Aggregate cohort characteristics are given in Table~\ref{tab:cohort}. De-identified per-patient evaluation tables (predicted and reference organism sets for \toolname, each conventional read-out and each directly prompted model) will be \todo{provided as Supplementary Data / made available on reasonable request}. PubMed identifiers and stored article judgements from the literature-retrieval and patient-matching analyses will be provided as Supplementary Data. Source data for Figs.~\ref{fig:benchmark} and \ref{fig:llm} are provided with this paper.

\subsection*{Code availability}
The \toolname pipeline, including all rules, prompts, configuration files, synthetic example data and the offline test suite, is available at \url{https://github.com/OliveWoo/Multimodal-Diagnosis-Model}; the version used for this manuscript is commit d82c0f7. Reproducing the language-model stages requires OpenAI API access and PubMed access; the deterministic stages and the replay of the reported decisions do not. The code will be released under \todo{licence}, with the version used here archived at Zenodo (\todo{DOI}). The scripts that generate the figures and tables of this manuscript are available at \url{https://github.com/ythuang0522/aetia-paper}.
